\RequirePackage{mymoniker}    % Replace CarrollCD in mymoniker.sty with your LastnameFirstNameMiddleName
\RequirePackage{dissertation} % required extra stuff
\documentclass[11pt]{article}
\usepackage[top=1.25in,bottom=1.25in,left=1.25in,right=1.25in]{geometry}
\usepackage{amsmath}

\pagestyle{empty}
\begin{document}


\begin{center}
{\Large {\bf Spiels for }}
\\[5pt]
{\large {\bf Your Name}}
\end{center}

\vspace{.3cm}

\section{\large Elevator Spiel} % This should be a 1 minute description of what you do that can be used when on an elevator ride, say

\section{\large Why do you want to work at this institution?}

There are 3 main reasons I would like to work at (your institution).

First, I enjoy thinking about challenging economic problems using mathematical modeling, theory, and computational tools. These are approaches that align well with the type of work that is conducted at (your institution).

Second, I'm excited to work with like-minded individuals whose diverse experiences and skills will complement my own. An environment like this will allow me to contribute to a greater, common goal.

Lastly, while I am finally finishing school, I do view myself as a lifelong learner. I'm eager to see how the concepts I've studied over the past decade translate into completing tasks in the work setting. At the very least, I feel the past 6 years of my life have helped me to hone my problem-solving skills.


\section{\large Job Market Paper Spiel}


\noindent
My job market paper is at the intersection of wealth inequality and macroeconomics. The core question I answer is: \textit{Does allowing for heterogeneous returns in a standard consumption-saving framework help explain the skewed distribution of wealth observed in the data?} 

\vspace{0.5em}
\noindent
The early literature on wealth inequality focused on whether skewed labor earnings alone could explain the shape of the wealth distribution. But while income is unequal, wealth is far more skewed, especially in the upper tail. Early models then tried adding stochastic returns to wealth as an exogenous process, which did help match the distribution empirically. However, those models don't relate wealth accumulation to  household behavior.

\vspace{0.5em}
\noindent
My paper is in line with heterogeneous agent macro models which contribute to this discussion. I show that introducing modest heterogeneity in returns on top of labor income risk significantly improves the model's ability to match targeted and untargeted wealth moments. Unlike time preferences or risk aversion, household returns are observable. The estimated distribution of returns from my paper corresponds to empirical evidence on returns from Norwegian administrative data. 

\vspace{0.5em}
\noindent
There are two key policy implications of my paper. First, capital income taxation reduce wealth inequality more than wealth taxation when returns differ across households. This is because 